\section{Introduction}

Near-Earth Objects (NEOs) are asteroids and comets whose orbit brings them close to the Earth.
NASA publishes their characteristics and their close approaches through the NeoWs API, but this
raw data is not organised for a non-specialist: finding the objects that will pass close to the
Earth next week, following a given asteroid over time or being warned when a large object comes
close all require some processing.

\vspace{0.25cm}

NEO-Watch is a web application that answers these needs. It imports the NASA data into its own
PostgreSQL database, exposes it through a REST API written with FastAPI, and offers a Streamlit
interface where each user has an account, a dashboard, favourite objects and personal alerts.

\vspace{0.25cm}

The analysis report, submitted in September, defined the scope, the features and the design of
the application. This final report describes what was actually built. Section~\ref{sec:specs}
compares the delivered features with the specifications. Section~\ref{sec:architecture} presents
the architecture and the main design choices, and section~\ref{sec:implementation} the
implementation of each feature. Section~\ref{sec:testing} describes the tests and the quality
tools, section~\ref{sec:organization} the organisation of the team, and the conclusion reviews the
results, the limits and the possible improvements. The individual assessments are in appendix.
